\documentclass[11pt,a4paper]{article}

\usepackage[utf8]{inputenc}
\usepackage[T1]{fontenc}
\usepackage{lmodern}
\usepackage{amsmath, amssymb, amsthm}
\usepackage{mathtools}
\usepackage{geometry}
\usepackage{hyperref}
\usepackage{xcolor}
\usepackage{booktabs}
\usepackage{enumitem}
\usepackage{microtype}
\usepackage{seqsplit}

\geometry{margin=2.8cm}
\setlength{\emergencystretch}{3em}
\sloppy
\newcommand{\ttbr}[1]{\texttt{\seqsplit{#1}}}
\hypersetup{
  colorlinks=true,
  linkcolor=blue!50!black,
  citecolor=blue!50!black,
  urlcolor=blue!50!black,
  pdftitle={Cross-AI, Lean-Verified Mathematics: A Case Study on the Collatz Conjecture (v5)},
  pdfauthor={Piero Borgatta}
}

\newtheorem{theorem}{Theorem}
\newtheorem{proposition}[theorem]{Proposition}
\newtheorem{lemma}[theorem]{Lemma}
\theoremstyle{definition}
\newtheorem{remark}[theorem]{Remark}

\newcommand{\Z}{\mathbb{Z}}
\newcommand{\Zp}{\mathbb{Z}_2}
\newcommand{\Q}{\mathbb{Q}}
\newcommand{\N}{\mathbb{N}}

\title{
  Cross-AI, Lean-Verified Mathematics: \\
  A Case Study on the Collatz Conjecture \\
  \large (v5 --- methodology retrospective and honest close)
}
\author{Piero Borgatta\thanks{Independent researcher.
\texttt{info@pieroborgatta.com}; ORCID 0009-0001-8025-2405.}}
\date{\today}

\begin{document}
\maketitle

\begin{abstract}
This is the fifth and closing version of a personal research program on the
Collatz conjecture. Its primary subject is the program's \emph{explicitly
declared secondary goal} from v2 onward: testing how far iterative,
cross-checked collaboration with large language model (LLM) assistants can
carry a non-standard attempt on an open mathematical problem. The
mathematics is reported honestly as a closed (but not abandoned) chapter,
and the methodology is brought to the foreground.

We report, with reproducible figures, what an AI pair (Anthropic Claude and
OpenAI Codex, with Google Gemini in the earliest phase) produced under the
author's direction over a five-week span: a mechanically verified,
\texttt{sorry}-/\texttt{admit}-/\texttt{axiom}-free Lean~4 + Mathlib core
(an exact congruential shadowing lemma, a no-infinite-shadowing corollary,
an expanding-cycle exclusion, finite Collatz--Wielandt spectral-radius
certificates, and an elementary descent bridge); a finite phantom-set
taxonomy; and roughly $1.4\times 10^{5}$ lines of AI-produced code and
prose, of which the author typed none by hand. We also report the
program's \emph{negative} results honestly --- the analytic routes that hit
documented barriers, including a cross-AI episode in which the assistants
generated, and then demolished, four ``revolutionary'' proposals --- and
we record the methodological observation we consider most interesting: a
cross-AI loop, when one side is anchored to concrete obstructions, can
\emph{recognize and prove the ceiling of its own approach}. We distill
transferable lessons and failure modes for AI-assisted theorem proving.
The mathematical frontier the program identified as the only productive one
(cycle exclusion) is left explicitly open as future work, not closed. We do
not claim a proof of the Collatz conjecture.
\end{abstract}

\tableofcontents

\section{The v1$\to$v5 arc}
\label{sec:arc}

This program began (v1) as an ambitious proposal: a spectral reduction of
the Collatz conjecture via ``phantom orbit shadowing,'' framed as a
\emph{possible route to a resolution}. It then evolved honestly:

\begin{itemize}[leftmargin=2em]
  \item \textbf{v1} --- the original spectral program, proposed as a possible
        Collatz resolution.
  \item \textbf{v2} --- the shadowing core mechanically verified in Lean~4;
        comparison with the concurrent work of \cite{chang2026}; the
        AI-collaboration methodology declared as an explicit secondary goal.
  \item \textbf{v3} --- a conditional spectral reduction, a phantom-set
        taxonomy at depth $K_0=16$, and an extended Lean operator/certificate
        layer.
  \item \textbf{v4} --- a consolidation that \emph{proved} the program's two
        analytic branches cannot reach the conjecture, mapping the two
        distinct barriers and isolating the open content precisely
        \cite{borgatta_v4}.
  \item \textbf{v5} (this document) --- the secondary goal becomes the
        subject: an honest methodology retrospective and the program's close,
        with the verified artifacts as evidence and the productive
        mathematical frontier left open.
\end{itemize}

We regard this arc --- from ``possible solution'' to a precise, honest map
of \emph{why} the approach stops and \emph{what} is nonetheless verified ---
as the scientifically valuable shape of the work, more so than any single
intermediate claim. Reporting it candidly, including the failures, is the
point of v5.

\section{The verified artifact}
\label{sec:artifact}

All of the following is formalized in Lean~4 + Mathlib and is
\texttt{sorry}-, \texttt{admit}- and \texttt{axiom}-free (verified by
\texttt{grep} over the \texttt{CollatzShadowing} sources and by
\texttt{lake build}).

\begin{itemize}[leftmargin=2em]
  \item \ttbr{exact\_shadowing}, \ttbr{exact\_shadowing\_periods} --- the
        exact congruential shadowing lemma.
  \item \ttbr{no\_infinite\_period\_congruence\_expansive} --- no infinite
        periodic shadowing.
  \item \ttbr{no\_positive\_endpoint\_eventually\_periodic\_expansive\_congruence}
        --- no positive integer on an expanding cycle (post-prefix form).
  \item \ttbr{spectralRadius\_le\_of\_finiteCWCertificate} --- the generic
        weighted Collatz--Wielandt bound, bridged to Mathlib's real
        \texttt{spectralRadius}.
  \item \ttbr{t10j32HighBitTailSpectralRadiusBound\_97\_2000} (single node,
        $\le 97/2000$) and
        \ttbr{k16s16KDeterministicGeneratedSpectralRadiusBound}
        (deterministic $(K,b)$ bound $< 3/4$ at $K_0=16$).
  \item \ttbr{classicalCollatz\_of\_uniformStrictDescent\_provedBridge} ---
        the elementary descent bridge to the classical conjecture.
\end{itemize}

The full declaration map is in \texttt{lean/CollatzShadowing/THEOREM\_INDEX.md};
the build is reproducible from \texttt{lean/} via \texttt{lake exe cache get}
then \texttt{lake build}.

\section{Quantitative summary}
\label{sec:stats}

All figures are read-only measurements of the repository
(\texttt{git}, \texttt{wc}, \texttt{grep}); reproduction commands and the
full breakdown are in \texttt{notes/PROJECT\_STATISTICS.md}.

\paragraph{Authorship.} Every file in the repository --- all code and all
prose --- was produced by the AI assistants under the author's direction.
\textbf{The author typed none of it by hand}; his contribution is direction,
conceptual framing, decisions, editing-by-instruction, and all external
actions (publication, compilation, version control operations).

\begin{center}
\begin{tabular}{lrl}
\toprule
category & lines & content \\
\midrule
AI-authored source & $\approx 61{,}800$ & Lean proofs $7{,}805$ + Python $54{,}009$ \\
generator-emitted & $43{,}513$ & Lean finite certificates (script output) \\
prose & $34{,}899$ & working notes $28{,}716$ + LaTeX $6{,}183$ \\
\midrule
total (AI-produced) & $\approx 140{,}200$ & human-typed: $0$ \\
\bottomrule
\end{tabular}
\end{center}

\begin{itemize}[leftmargin=2em]
  \item \textbf{Lean:} 40 files; 302 AI-authored theorems/lemmas (14 modules)
        plus 966 emitted by generators (26 certificate modules); \emph{zero}
        \texttt{sorry}/\texttt{admit}/\texttt{axiom}.
  \item \textbf{Python:} 151 scripts (all tracked, none missed) ---
        spectral program 83, phase-one empirical 42, phantom taxonomy 16,
        root-level 10.
  \item \textbf{History:} 84 commits over 37 days ($\approx$ 5.3 weeks);
        commit authorship (via \texttt{.mailmap}) Codex 52, Piero Borgatta 35.
  \item \textbf{Honesty caveat:} never quote the 51k Lean total or the 1{,}268
        theorem count without the AI-authored vs generator-emitted split; the
        generator output dominates both.
\end{itemize}

\paragraph{Soft metrics.} Human hours, number of AI sessions, and message
volume are not cleanly measurable from the repository and are reported
separately as labelled estimates; we deliberately do \emph{not} state an
``AI work hours'' figure, since LLM inference time is not comparable to human
labour hours.

\section{Methodology: the cross-AI protocol}
\label{sec:method}

The working configuration was:
\begin{enumerate}[leftmargin=2em]
  \item \textbf{Human direction.} The author set the research direction and
        framing, made all strategic and publication decisions, and performed
        all external actions.
  \item \textbf{Generate / critique split.} The assistants alternated roles:
        one proposing constructions or proofs, the other auditing them
        against the concrete state of the problem.
  \item \textbf{Lean as ground truth.} Every formal claim was funnelled
        through Lean: a proof counts only if \texttt{lake build} succeeds
        with no \texttt{sorry}. This converts ``plausible-looking'' AI output
        into a binary, machine-checked verdict --- the single most important
        discipline of the project.
\end{enumerate}

The decisive methodological lesson (Section~\ref{sec:limits}) is that the
generate/critique split is productive \emph{only when the critic is anchored
to concrete objects and obstructions}, not to labels. The same kind of model,
in the generative role, repeatedly produced plausible but structurally empty
proposals; in the critic role, with the problem's barriers loaded as context,
it dismantled them. The asymmetry is one of \emph{role and discipline}, not of
model capability.

\section{What worked}
\label{sec:worked}

\begin{itemize}[leftmargin=2em]
  \item Producing and \emph{mechanically verifying} a real (if elementary)
        body of mathematics: the shadowing core and the finite certificates.
  \item Managing a large, multi-week formalization with a coherent module
        structure and a clean \texttt{sorry}-free state.
  \item Generating exact finite certificates programmatically and importing
        them into Lean with machine-checked row witnesses.
  \item Cross-AI auditing that repeatedly caught Mathlib API mismatches and
        Lean~3-vs-4 syntax errors one assistant alone missed.
\end{itemize}

\section{What failed --- honest failure modes}
\label{sec:failed}

This is, in our view, the most useful part of the report: most accounts of
``LLMs doing mathematics'' present only successes.

\begin{itemize}[leftmargin=2em]
  \item \textbf{Abandoned analytic branches.} A $\Zp$ martingale Banach space
        (failed: negative 2-adic child-cylinder signal); an A0 weak averaged
        bridge and an A1 refined-square kernel (partial; barrier); a pointwise
        first-barrier branch (conditional; negatively resolved in v4).
  \item \textbf{``Prestige-transfer'' proposals.} In a documented cross-AI
        episode, the assistants generated four ``revolutionary'' pivots
        (Galois/anabelian action on the inverse tree; non-standard models /
        independence; ``exotic'' Banach spaces; holomorphic extension), each
        of which was then dismantled with a specific structural reason. The
        common failure: each imports rigidity from a richer ambient category
        (algebraic, logical, measure-analytic, complex-analytic), but the hard
        invariant --- an individual integer carrying an aperiodic $2$-adic
        valuation word --- is precisely what every such embedding renders
        generic or invisible.
  \item \textbf{Label pattern-matching.} Concrete instances: proposing
        Mathlib contributions based on \emph{removed} API
        (\texttt{Nat.bit0}/\texttt{bit1}); hallucinated lemma names; assuming
        novel general-purpose lemmas where the project had in fact
        \emph{assembled} existing Mathlib primitives.
  \item \textbf{Generative overreach.} A persistent tendency to connect
        semantically adjacent domains without checking structural
        compatibility against the actual objects.
\end{itemize}

\section{The limits, recognized and proved}
\label{sec:limits}

The program's strongest methodological outcome is that the cross-AI loop did
not merely fail quietly; it \emph{localized and proved its own ceiling}.

\begin{itemize}[leftmargin=2em]
  \item \textbf{Two barriers, separated.} v4 establishes that Collatz splits
        into two distinct hard halves: \emph{no divergent orbits}
        (barrier: distributional-to-pointwise) and \emph{no nontrivial cycles}
        (barrier: Diophantine approximation of $\log_2 3$). The program's
        entire analytic effort lived in the first.
  \item \textbf{A negative structural result.} Within the phantom/$2$-adic
        formalism, the rigidity one would need on the orbit limit point is
        \emph{logically equivalent} to the divergence conjecture itself
        (a non-transport observation: the parity-vector conjugacy carries
        eventual periodicity to dynamical pre-periodicity, not to
        rationality; conditional on the standard Lagarias coding).
\end{itemize}

That an AI-assisted process can arrive at, and rigorously articulate, ``this
approach cannot work, and here is why'' --- rather than continuing to produce
hopeful variations --- is a result we did not anticipate and consider the
central methodological finding.

\section{Open mathematical direction (left open, not closed)}
\label{sec:open}

The program is closed as a methodology study, but the mathematics it pointed
to is \emph{not abandoned}. v4 identified one genuinely productive frontier,
and we keep it explicitly open as future work:

\begin{quote}
The \textbf{cycle half} (no nontrivial cycles). This is the only half whose
tools (linear forms in logarithms; continued fractions of $\log_2 3$;
computational verification; \cite{steiner1977, simons2005, simonsdeweger2005,
barina2021}) genuinely bite. The realistic, finishable target is the
\emph{structural layer} formalized in Lean --- the cycle equation
$n\,(2^A - 3^L) = C_w$, the contractivity condition $2^A > 3^L$ for positive
cycles, the Diophantine condition on $A/L$, and the link to computational
verification --- with the deep analytic exclusions left as named hypotheses,
since a usable form of Baker's theorem is not currently in Mathlib.
\end{quote}

Anyone (the author included) may resume the program from this point. The
verified Lean library (Section~\ref{sec:artifact}) is the natural foundation
for that continuation.

\section{Lessons for AI-assisted mathematics}
\label{sec:lessons}

\begin{enumerate}[leftmargin=2em]
  \item \textbf{Mechanical verification is the load-bearing discipline.}
        Lean's binary verdict is what separates this program's verified core
        from the unverifiable analytic speculation; AI prose proofs without it
        are not evidence.
  \item \textbf{Cross-AI value is adversarial, not additive.} The benefit came
        from one model auditing another against concrete obstructions, not
        from two models agreeing.
  \item \textbf{Generative fluency is not mathematical judgment.} The same
        system produces both verified proofs and structurally empty
        ``revolutionary'' proposals; the discriminator is anchoring to objects,
        not to vocabulary.
  \item \textbf{Honest negative results are a deliverable.} Localizing and
        proving an approach's ceiling saved further wasted effort and is, we
        argue, the most transferable output here.
  \item \textbf{Beware inflated metrics.} Report AI-authored vs generated
        separately; do not fabricate ``AI hours.''
\end{enumerate}

\section{Reproducibility}
\label{sec:repro}

The Lean~4 library, the Python scripts, the working notes (including
\texttt{PROJECT\_STATISTICS.md} with reproduction commands), and the paper
sources are in the supplementary archive and at
\url{https://github.com/PieroBorgatta/Collatz}. Earlier versions: concept DOI
\texttt{10.5281/zenodo.20021537}; v4 \texttt{10.5281/zenodo.20544464}.

\section{Conclusion}
\label{sec:conclusion}

Over five weeks, an AI pair under human direction produced a mechanically
verified mathematical core, a precise map of two barriers blocking the
Collatz conjecture, an honest catalogue of failed routes, and --- the result
we find most interesting --- a proof of its own approach's ceiling. We do not
claim progress on the conjecture. We claim a verified artifact, a candid
account of what cross-AI mathematical collaboration can and cannot do, and a
clearly delimited open frontier for whoever continues.

\begin{thebibliography}{9}

\bibitem{borgatta_v4}
P.~Borgatta, \emph{Phantom Orbit Shadowing for the Collatz Conjecture:
Verified Core, Two Barriers, and a Repositioning} (v4), Zenodo, DOI
\texttt{10.5281/zenodo.20544464}, 2026.

\bibitem{chang2026}
(Chang), \emph{Exploring Collatz Dynamics with Human--LLM Collaboration},
2026 (concurrent, independent; cited as in v2--v4).

\bibitem{steiner1977}
R.~P.~Steiner, \emph{A theorem on the Syracuse problem}, Proc. 7th Manitoba
Conf. on Numerical Mathematics, 1977, 553--559.

\bibitem{simons2005}
J.~L.~Simons, \emph{On the nonexistence of $2$-cycles for the $3x+1$ problem},
Math. Comp. \textbf{74} (2005), 1565--1572.

\bibitem{simonsdeweger2005}
J.~L.~Simons and B.~M.~M.~de Weger, \emph{Theoretical and computational bounds
for $m$-cycles of the $3n+1$ problem}, Acta Arith. \textbf{117} (2005), 51--70.

\bibitem{barina2021}
D.~Barina, \emph{Convergence verification of the Collatz problem},
J. Supercomput. \textbf{77} (2021), 2681--2688.

\end{thebibliography}

\end{document}
